前面的 Pauli--Lubanski 算符已经告诉我们，质量和自旋可以作为单粒子态的 Lorentz 不变量标签。但还缺少一个更具体的问题：把粒子从一个惯性系变到另一个惯性系时，四动量会改变，内部自旋指标究竟怎样跟着变？

固定一条质量壳后，任意四动量 $p$ 都可由某个标准 Lorentz 变换从参考四动量 $\bar p$ 得到。若再施加一般 Lorentz 变换 $\Lambda$，把 $\Lambda p$ 重新拉回 $\bar p$ 所留下的复合变换会保持 $\bar p$ 不变，因此它不能改变参考动量，只能作用在内部指标上。所有保持 $\bar p$ 不变的 Lorentz 变换组成\emph{小群}。有质量粒子的自旋和无质量粒子的螺旋度，正是这个剩余对称群的不可约表示标签。


为使 Lorentz 变换与单粒子完备关系兼容，单粒子动量本征态采用协变归一化。对内部量子数 $r$，取
\begin{equation}
 \boxed{
 \langle \mathbf p',r'|\mathbf p,r\rangle_{\rm cov}
 =2p^0(2\pi\hbar)^3\delta_{r'r}
 \delta^{(3)}(\mathbf p'-\mathbf p),
 \qquad p^0=\frac{E_{\mathbf p}}{c},}
 \label{eq:rel-state-normalization-si-r9}
\end{equation}
并配套使用
\begin{equation}
 \boxed{
 \id(1)=\sum_r\int
 \frac{\dd^3 p}{(2\pi\hbar)^3\,2p^0}
 |\mathbf p,r\rangle_{\rm cov}
 {}_{\rm cov}\langle\mathbf p,r|.}
 \label{eq:rel-completeness-si-r9}
\end{equation}
这组归一化使测度 $\dd^3 p/(2p^0)$ 保持 Lorentz 不变，因此 Wigner 旋转的内部矩阵可以与动量测度分离，不会额外出现由非协变三维 $\delta$ 归一化产生的 Jacobian 平方根。

固定一条给定质量壳，并在这条轨道上选取标准动量 $\bar p^\mu$。对每个同轨道四动量 $p^\mu$，选一个标准 Lorentz 变换 $L(p)$ 满足
\begin{equation}
 L(p)\bar p=p.
 \label{eq:wigner-standard-boost-dp9}
\end{equation}
以 $|\bar p,\sigma\rangle$ 表示标准动量上的的内部基，定义
\begin{equation}
 |p,\sigma\rangle
 \equiv U[L(p)]|\bar p,\sigma\rangle.
 \label{eq:wigner-state-definition-dp9}
\end{equation}
$\sigma$ 标记固定标准动量时仍可变化的内部自由度，不增加时空坐标。对任意 Lorentz 变换 $\Lambda$，定义
\begin{equation}
 \boxed{
 W(\Lambda,p)
 \equiv L^{-1}(\Lambda p)\,\Lambda\,L(p),
 \qquad
 W(\Lambda,p)\bar p=\bar p.}
 \label{eq:wigner-little-group-element-dp68}
\end{equation}
因此 $W(\Lambda,p)$ 不再改变四动量，只能作用在内部指标上。若内部空间采用小群不可约表示 $D$，单粒子态的完整 Lorentz 变换为
\begin{equation}
 \boxed{
 U(\Lambda)|p,\sigma\rangle
 =\sum_{\sigma'}
 D_{\sigma'\sigma}[W(\Lambda,p)]
 |\Lambda p,\sigma'\rangle.}
 \label{eq:wigner-state-transform-dp68}
\end{equation}

\begin{figure}[htbp]
 \centering
 \resizebox{0.92\textwidth}{!}{\begin{tikzpicture}[x=1cm,y=1cm,>=Stealth,
 every node/.style={font=\large},
 box/.style={draw,rounded corners=2pt,inner xsep=8pt,inner ysep=5pt,align=center,line width=.75pt},
 line label/.style={fill=white,fill opacity=.96,text opacity=1,inner sep=1.2pt}]
\node[box] (k) at (-4.5,0) {standard momentum\\$\bar p$};
\node[box] (p) at (-1.2,2.0) {$p=L(p)\bar p$};
\node[box] (lp) at (3.1,2.0) {$\Lambda p$};
\node[box] (kb) at (3.1,-1.0) {$\bar p$\\same orbit reference};

\draw[->,line width=.85pt] (k)--node[line label,above left] {$L(p)$}(p);
\draw[->,line width=.85pt] (p)--node[line label,above] {$\Lambda$}(lp);
\draw[->,line width=.85pt] (lp)--node[line label,right] {$L^{-1}(\Lambda p)$}(kb);
\draw[->,line width=.95pt,dashed] (k)--node[line label,below] {$W(\Lambda,p)$}(kb);

\node[box,align=center,font=\normalsize] (stab) at (-0.7,-2.65)
  {$G_{\bar p}=\{\Lambda\mid \Lambda\bar p=\bar p\}$\\stabilizer of the reference momentum};
\node[box] (massive) at (-3.2,-4.65) {$\bar p^2=m^2c^2>0$\\$G_{\bar p}\simeq SO(3)$\\double cover $SU(2)$\\spin $s$};
\node[box] (massless) at (3.4,-4.65) {$\bar k^2=0$\\$G_{\bar k}\simeq ISO(2)$\\discrete-helicity sector\\helicity $\lambda$};
\draw[->,line width=.8pt] (stab.south west) --node[line label,left,font=\normalsize]{timelike orbit}(massive.north);
\draw[->,line width=.8pt] (stab.south east) --node[line label,right,font=\normalsize]{null orbit}(massless.north);
\end{tikzpicture}
}
 \caption{Wigner 小群构造。$L(p)$ 把标准动量推到 $p$，而 $W(\Lambda,p)=L^{-1}(\Lambda p)\Lambda L(p)$ 把复合变换拉回同一个标准动量，因此只作用在内部自旋/螺旋度空间。}
 \label{fig:wigner-little-group-dp9}
\end{figure}

对 $m>0$，最方便的标准动量就是静止四动量 $\bar p_m^\mu=(mc,\mathbf0)$。保持它不变的 Lorentz 变换不能混合时间与空间分量，只能在三维空间做转动，因此
\begin{equation}
 \boxed{
 G_{\bar p_m}\simeq SO(3),
 \qquad
 \widetilde G_{\bar p_m}\simeq SU(2).}
 \label{eq:wigner-massive-little-group-dp68}
\end{equation}
量子态使用双覆盖 $SU(2)$，所以其不可约表示由 $s=0,\tfrac12,1,\ldots$ 标记，并具有 $2s+1$ 个自旋态。这与前文 $W^2=-m^2c^2\hbar^2s(s+1)$ 的 Casimir 分类是同一件事的两种表述。

无质量轨道没有静止系，因此不能把有质量构造简单令 $m\to0$。选沿 $+z$ 传播的标准零质量动量
\begin{equation}
 \bar k^\mu=(\kappa_\star,0,0,\kappa_\star),
 \qquad \kappa_\star>0,
 \qquad \bar k^2=0,
 \label{eq:wigner-massless-standard-momentum-dp9}
\end{equation}
其中 $\kappa_\star$ 具有物理动量量纲，其数值只是在零质量轨道上选参考点。全部稳定子方向由满足 $X\bar k=0$ 的 Lorentz Lie 代数元给出。写一般元
\begin{equation*}
 X=\sum_{i=1}^3(a_iJ_i+b_iK_i)
\end{equation*}
出发，并要求它在 $P^\mu=\bar k^\mu$ 上满足 $[X,P^\mu]=0$。四个分量的线性条件化为
\begin{equation*}
 b_3=0,
 \qquad
 b_1=-a_2,
 \qquad
 b_2=a_1,
\end{equation*}
所以稳定子 Lie 代数恰为三维。取一组方便的基为
\begin{equation}
 \mathsf N_1\equiv K_1-J_2,
 \qquad
 \mathsf N_2\equiv K_2+J_1,
 \qquad
 J_3.
 \label{eq:wigner-massless-little-group-generators-dp68}
\end{equation}
它们满足
\begin{equation}
 \boxed{
 [J_3,\mathsf N_1]=\ii\hbar\mathsf N_2,
 \qquad
 [J_3,\mathsf N_2]=-\ii\hbar\mathsf N_1,
 \qquad
 [\mathsf N_1,\mathsf N_2]=0,}
 \label{eq:wigner-massless-little-group-algebra-dp68}
\end{equation}
这正是二维平面转动加平移的代数，因此在 $SO^+(1,3)$ 中零质量标准动量的小群同构于 $ISO(2)$。这里的“平移”作用在内部表示空间，不是把粒子的时空位置平移。量子态若由 Lorentz 群的双覆盖 $SL(2,\mathbb C)$ 实现，则实际作用的是 $ISO(2)$ 的相应双覆盖；这一区分对半整数螺旋度是必要的。

这个小群的量子不可约表示有两类。若 $\mathsf N_{1,2}$ 非平凡作用，会得到连续自旋表示；若这两个“平移型”生成元在物理态上平凡作用，则只剩绕运动方向的转动生成元 $J_3$。这种第二类表示由一个离散或半整数的螺旋度标签 $\lambda$ 区分，并满足
\begin{equation}
 \mathsf N_1|\bar k,\lambda\rangle=0,
 \qquad
 \mathsf N_2|\bar k,\lambda\rangle=0,
 \qquad
 J_3|\bar k,\lambda\rangle=\hbar\lambda|\bar k,\lambda\rangle.
 \label{eq:massless-helicity-irrep-dp9}
\end{equation}
此时 Pauli--Lubanski 矢量与四动量平行：
\begin{equation}
 \boxed{W^\mu|k,\lambda\rangle
 =\hbar\lambda\,P^\mu|k,\lambda\rangle.}
 \label{eq:massless-pauli-lubanski-helicity-dp9}
\end{equation}
因此无质量表示的 $W^2=0$ 对所有离散螺旋度都相同；真正区分内部态的是小群转动生成元本征值 $\lambda$。光子的物理 Hilbert 空间取 $\lambda=\pm1$，无质量自旋 $1/2$ 场取 $\lambda=\pm\tfrac12$。光子的 $\lambda=0$ 不属于物理横向单光子空间；后面 Maxwell 约束与规范商空间会从场变量层面再次得到同一结论。

三个常被混写的对象属于不同的表示层级：$(j_L,j_R)$ 是\emph{局域场}在有限维 Lorentz 复表示中的变换类型；$s$ 是有质量单粒子小群双覆盖 $SU(2)$ 的不可约标签；$\lambda$ 是无质量离散表示中稳定子转动方向的本征值。对整数自旋可直接把它看成 $SO(2)$ 角色标签；对半整数自旋必须在 Lorentz 双覆盖中的小群上理解它。一个四矢量场属于 $(\tfrac12,\tfrac12)$ 并不意味着物理光子有四个自旋态；规范约束最终只留下 $\lambda=\pm1$ 两个单粒子螺旋度态。


\begin{derivation}[从\eqnrefs{eq:wigner-state-definition-dp9}到\eqnrefs{eq:wigner-state-transform-dp68}]
对任意保向正时向 Lorentz 变换 $\Lambda$，定义

\begin{equation*}
 {
 W(\Lambda,p)
 \equiv L^{-1}(\Lambda p)\,\Lambda\,L(p).}
 \end{equation*}
直接作用在标准动量上：
\begin{align*}
 W(\Lambda,p)\bar p
 &=L^{-1}(\Lambda p)\Lambda L(p)\bar p,\\
 &=L^{-1}(\Lambda p)\Lambda p,\\
 &=\bar p.
 \end{align*}
因此 $W(\Lambda,p)$ 确实属于稳定 $\bar p$ 的小群
\begin{equation*}
 G_{\bar p}\equiv\{W\in SO^+(1,3)\mid W\bar p=\bar p\}.
 \end{equation*}
再把 \eqnrefs{eq:wigner-state-definition-dp9} 代入态变换：
\begin{align*}
 U(\Lambda)|p,\sigma\rangle
 &=U(\Lambda)U[L(p)]|\bar p,\sigma\rangle,\\
 &=U[L(\Lambda p)]U[W(\Lambda,p)]|\bar p,\sigma\rangle.
\end{align*}
因为 $W$ 不改变 $\bar p$，它只能在固定动量的内部空间中作用。若该内部空间取小群不可约表示 $D$，
\begin{equation*}
 U(W)|\bar p,\sigma\rangle
 =\sum_{\sigma'}D_{\sigma'\sigma}(W)|\bar p,\sigma'\rangle,
\end{equation*}
于是
\begin{equation*}
 {
 U(\Lambda)|p,\sigma\rangle
 =\sum_{\sigma'}
 D_{\sigma'\sigma}[W(\Lambda,p)]
 |\Lambda p,\sigma'\rangle.}
 \end{equation*}
这条式子把四动量的 Lorentz 变换与内部自旋/螺旋度变换严格分离。若改用非协变三维 $\delta$ 归一化，右端还会出现由雅可比行列式产生的 $\sqrt{(\Lambda p)^0/p^0}$ 因子；后续场量子化采用协变归一化，因此不再引入这一额外因子。
\end{derivation}
\begin{derivation}[从\eqnrefs{eq:wigner-little-group-element-dp68}到\eqnrefs{eq:wigner-massive-little-group-dp68}]
有质量粒子可选静止标准四动量
\begin{equation*}
 \bar p_m^\mu=(mc,0,0,0).
\end{equation*}
小群元素 $W$ 的定义是 $W\bar p_m=\bar p_m$。逐分量写出
\begin{equation*}
 W^\mu{}_{\nu}\bar p_m^\nu
 =mc\,W^\mu{}_0
 =mc\,\delta^\mu{}_0,
\end{equation*}
所以第一列被完全固定：
\begin{equation*}
 W^0{}_0=1,
 \qquad
 W^i{}_0=0.
\end{equation*}
还需要证明第一行也没有 boost 分量。Lorentz 条件
$W^{\mathsf T}\eta W=\eta$ 的 $(0,i)$ 分量给
\begin{equation*}
 \eta_{\alpha\beta}W^\alpha{}_0W^\beta{}_i=0.
\end{equation*}
代入 $W^\alpha{}_0=\delta^\alpha{}_0$ 后得到
$W^0{}_i=0$。因此 $W$ 必须具有块形式
\begin{equation*}
 W=
 \begin{pmatrix}
 1&\mathbf0^{\mathsf T}\\
 \mathbf0&R
 \end{pmatrix}.
\end{equation*}
再把这个块矩阵代回 Lorentz 条件，空间块满足
\begin{equation*}
 R^{\mathsf T}R=I_3.
\end{equation*}
若只考虑与恒等变换连续相连的 proper orthochronous Lorentz 群，还要求
$\det W=+1$，因此 $\det R=+1$。所以有质量标准动量的连续稳定子正是
$SO(3)$。

量子态的自旋不是用 $SO(3)$ 的普通表示分类，而是用其双覆盖 $SU(2)$ 的幺正不可约表示分类，因为 $2\pi$ 空间转动可以在态矢量上产生负号而不改变物理射线。$SU(2)$ 不可约表示由
$s=0,\tfrac12,1,\ldots$ 标记，维数为 $2s+1$；这正是正文\eqnrefs{eq:wigner-massive-little-group-dp68} 中出现的自旋量子数，也与 Pauli--Lubanski Casimir 的标签一致。
\end{derivation}
\begin{derivation}[从\eqnrefs{eq:wigner-massless-standard-momentum-dp9}到\eqnrefs{eq:wigner-massless-little-group-algebra-dp68}]

稳定子代数的完备性由一般 Lorentz 代数元的线性稳定条件决定。先解出全部稳定方向，再计算所得三维解空间的 Lie 括号。

写
\begin{equation*}
 X=a_1J_1+a_2J_2+a_3J_3
 +b_1K_1+b_2K_2+b_3K_3.
\end{equation*}
由正文\eqnrefs{eq:poincare-mp}，
\begin{equation*}
 [J_i,P^0]=0,
 \qquad
 [J_i,P^j]=\ii\hbar\epsilon_{ijk}P^k,
\end{equation*}
以及
\begin{equation*}
 [K_i,P^0]=-\ii\hbar P^i,
 \qquad
 [K_i,P^j]=-\ii\hbar\delta_{ij}P^0.
\end{equation*}
在
\begin{equation*}
 \bar k^\mu=(\kappa_\star,0,0,\kappa_\star)
\end{equation*}
上逐分量要求 $[X,P^\mu]_{P=\bar k}=0$。时间分量给
\begin{equation*}
 -\ii\hbar\kappa_\star b_3=0
 \quad\Longrightarrow\quad b_3=0.
\end{equation*}
第一空间分量给
\begin{equation*}
 -\ii\hbar\kappa_\star(a_2+b_1)=0
 \quad\Longrightarrow\quad b_1=-a_2,
\end{equation*}
第二空间分量给
\begin{equation*}
 \ii\hbar\kappa_\star(a_1-b_2)=0
 \quad\Longrightarrow\quad b_2=a_1.
\end{equation*}
第三空间分量在这些条件下自动为零。因此任意稳定子元都唯一写成
\begin{equation*}
 X=a_3J_3+a_1(K_2+J_1)-a_2(K_1-J_2),
\end{equation*}
证明稳定子空间恰为
\begin{equation*}
 \operatorname{span}\{J_3,\mathsf N_1,\mathsf N_2\}.
\end{equation*}

再用正文\eqnrefs{eq:JK-algebra}：
\begin{align*}
 [J_3,\mathsf N_1]
 &=[J_3,K_1]-[J_3,J_2]
 =\ii\hbar(K_2+J_1)
 =\ii\hbar\mathsf N_2,\\
 [J_3,\mathsf N_2]
 &=[J_3,K_2]+[J_3,J_1]
 =-\ii\hbar(K_1-J_2)
 =-\ii\hbar\mathsf N_1,\\
 [\mathsf N_1,\mathsf N_2]
 &=[K_1,K_2]+[K_1,J_1]-[J_2,K_2]-[J_2,J_1]\\
 &=-\ii\hbar J_3+\ii\hbar J_3\\
 &=0.
\end{align*}
这就是正文\eqnrefs{eq:wigner-massless-little-group-algebra-dp68}。由于三维解空间已经由一般稳定条件完整求出，所得 $\mathfrak{iso}(2)$ 不再只是三个候选生成元的子代数，而是零质量标准动量的全部连续稳定子 Lie 代数。
\end{derivation}
